% ============================================================
%  main-groupmeeting.tex — 组会汇报 Beamer 演示文稿模板
%  ============================================================
%  使用方式：
%    1. 修改下方「元数据」区的标题、汇报人、导师等信息
%    2. 在各帧中替换示例内容为你的本周工作
%    3. 将图片放入 img/ 目录，用 \includegraphics{img/xxx.png} 引用
%    4. 终端运行以下命令编译：
%       latexmk -jobname=main-groupmeeting main-groupmeeting.tex
%    5. 查看生成的 main-groupmeeting.pdf
%  ============================================================

\documentclass[aspectratio=169]{beamer}

% ---------- Style ----------
\usepackage{beamer-style}

% ============================================================
%  元数据 — 修改为你自己的信息
% ============================================================
\title{本周研究进展汇报}
\subtitle{课题组组会}
\author{张\textbf{三}}
\institute{某某大学 ~ 计算机科学与技术学院 ~ 某某课题组}
\date{2026 年 6 月 19 日}

\begin{document}

% ============================================================
%  1. 标题页 (Title Page)
% ============================================================

\begin{frame}[shrink]
  \titlepage

  \vspace{-1.5em}
  \begin{center}
    \footnotesize
    \textbf{指导教师：}李四 ~ 教授
  \end{center}
\end{frame}

% ============================================================
%  2. 本周工作概览 (Weekly Overview)
% ============================================================

\begin{frame}{本周工作概览}
  \begin{itemize}
    \item \textbf{工作一：}完成了实验数据采集与预处理，共收集 500 条样本。
    \item \textbf{工作二：}实现了基线模型并完成初步训练，准确率达到 85.3\%。
    \item \textbf{工作三：}阅读了 2 篇相关论文，整理了方法对比笔记。
  \end{itemize}

  \vspace{1em}
  \begin{block}{总体进度}
    按计划推进，实验部分比预期提前 1 天完成。
  \end{block}
\end{frame}

% ============================================================
%  3. 工作进展详情 (Detailed Progress)
%     可按需增减帧数，每帧展开一项工作
% ============================================================

\begin{frame}{工作进展详情 — 数据预处理}
  \begin{itemize}
    \item 数据来源：公开数据集 + 自建标注数据（共 500 条）。
    \item 预处理步骤：
      \begin{enumerate}
        \item 文本清洗（去除噪声、统一编码为 UTF-8）。
        \item 分词与词性标注（使用 jieba 分词器）。
        \item 构建词表并过滤低频词（频次 < 3 的词语被替换为 \texttt{<UNK>}）。
      \end{enumerate}
    \item 最终构建了包含 12,000 个词语的词表。
  \end{itemize}

  \begin{exampleblock}{经验记录}
    jieba 分词在中文科技文本上效果良好，但部分专业术语需要手动添加自定义词典。
  \end{exampleblock}
\end{frame}

\begin{frame}{工作进展详情 — 基线模型实验}
  \begin{columns}[T]
    \begin{column}{0.48\textwidth}
      \textbf{实验设置}
      \begin{itemize}
        \item 模型：BERT-base-Chinese
        \item 训练集：400 条
        \item 验证集：50 条
        \item 测试集：50 条
        \item Epochs：10
        \item Batch size：16
      \end{itemize}
    \end{column}
    \begin{column}{0.48\textwidth}
      \textbf{实验结果}
      \begin{itemize}
        \item 验证集准确率：\textbf{85.3\%}
        \item 测试集准确率：\textbf{84.1\%}
        \item 训练耗时：约 12 分钟 (GPU)
      \end{itemize}

      \vspace{1em}
      % 如果你的实验有图表，替换为实际图片：
      % \includegraphics[width=\columnwidth]{img/your-result.png}
    \end{column}
  \end{columns}
\end{frame}

% ============================================================
%  4. 文献阅读（可选帧）
%     如果本周读了论文，取消注释使用此帧
% ============================================================

% \begin{frame}{文献阅读笔记}
%   \begin{block}{论文标题一 [arXiv 2026]}
%     \begin{itemize}
%       \item 核心方法：提出了一种基于注意力机制的改进方案。
%       \item 关键结论：在 X 数据集上提升了 3.2 个百分点的准确率。
%       \item 与本文关系：可以尝试将其方法迁移到我们的任务中。
%     \end{itemize}
%   \end{block}
%
%   \vspace{0.5em}
%   \begin{block}{论文标题二 [ACL 2025]}
%     \begin{itemize}
%       \item 核心方法：使用对比学习增强文本表示。
%       \item 关键结论：在小样本场景下效果显著。
%       \item 与本文关系：我们的数据量有限，可参考其小样本策略。
%     \end{itemize}
%   \end{block}
% \end{frame}

% ============================================================
%  5. 遇到的问题与讨论 (Problems & Discussion)
%     组会的核心环节：把困难说出来，获得导师和同门的建议
% ============================================================

\begin{frame}{遇到的问题与讨论}
  \begin{alertblock}{问题一：模型在小样本类别上效果差}
    \begin{itemize}
      \item 现象：测试集中样本数 < 10 的类别，F1 值仅 0.42。
      \item 已尝试：数据增强（回译、同义词替换）→ 提升不明显。
      \item 待讨论：是否可以合并低频类别？或引入外部知识图谱？
    \end{itemize}
  \end{alertblock}

  \vspace{0.5em}
  \begin{alertblock}{问题二：训练时间超出预期}
    \begin{itemize}
      \item 现象：全量数据训练一次约需 45 分钟（CPU），GPU 资源紧张。
      \item 已尝试：减小 batch size → 收敛变慢。
      \item 待讨论：是否可以使用混合精度训练？或有其他 GPU 资源可用？
    \end{itemize}
  \end{alertblock}
\end{frame}

% ============================================================
%  6. 下周工作计划 (Next Week Plan)
% ============================================================

\begin{frame}{下周工作计划}
  \begin{enumerate}
    \item \textbf{模型改进}（预计 3 天）
      \begin{itemize}
        \item 调研注意力机制变体方案。
        \item 实现 2 种改进模型并对比效果。
      \end{itemize}
    \item \textbf{实验补充}（预计 2 天）
      \begin{itemize}
        \item 补充消融实验（验证各模块的贡献）。
        \item 整理实验数据表格和图表。
      \end{itemize}
    \item \textbf{论文撰写}（预计 2 天）
      \begin{itemize}
        \item 更新 Related Work 部分（新增 2 篇引用）。
        \item 完成实验分析部分的初稿。
      \end{itemize}
  \end{enumerate}

  \vspace{0.5em}
  \begin{block}{目标}
    下周末前完成改进模型的初步实验，并提交实验分析草稿给导师审阅。
  \end{block}
\end{frame}

% ============================================================
%  7. 致谢 & 讨论 (Thank You & Q&A)
% ============================================================

\begin{frame}[c]{}
  \centering
  \Huge 谢谢！\\[0.5em]
  \Large 请各位老师和同学批评指正
\end{frame}

% ============================================================
%  8. 附录 (Appendix)
%     附录帧页码自动显示为 A-1, A-2, ...，
%     适合放置备份材料供讨论环节按需展示。
%     删除或替换为你的实际附录内容。
% ============================================================

\appendix

\section{附录}

\begin{frame}{附录 A：补充材料}
  \begin{itemize}
    \item 附录内容可包括：详细实验数据、代码片段、额外图表、文献阅读笔记详情等。
    \item 附录幻灯片页码自动显示为 \texttt{A-1, A-2, ...}，与正文 \texttt{X/Y} 格式区分。
    \item 根据需要可添加更多附录帧，页码自动递增。
  \end{itemize}

  \vspace{1em}
  \begin{block}{使用提示}
    汇报时不必逐页放映附录；当讨论涉及细节时，快速跳转至对应幻灯片展示。
  \end{block}
\end{frame}

% --- 如有需要，取消注释以下帧以在附录中列出参考文献 ---

% \begin{frame}[allowframebreaks]{参考文献}
%   \nocite{*}
%   \printbibliography
% \end{frame}

\end{document}
